\section{与 Parag 等信息论框架的深度比较}
\label{sec:part14}

\subsection{Parag 等经典泊松更新信息论回顾}

近年来，牛津大学的 Kris Parag 与 Christl Donnelly \citep{parag2022quantifying, parag2022fundamental} 将经典信息论与统计物理引入传染病再生数估计，建立了以 Fisher 信息量与互信息（Mutual Information）为核心的评估框架。其基准模型设定如下：
\begin{enumerate}
  \item 发病计数服从泊松分布：$I_t \mid \Lambda_t, R_t \sim \text{Poisson}(R_t \Lambda_t)$；
  \item 其条件对数似然为 $\ell(R_t) = I_t \ln(R_t \Lambda_t) - R_t \Lambda_t - \ln(I_t!)$；
  \item 经典泊松 Fisher 信息量为：
  \begin{equation}
    \mathcal{I}_{\text{Poisson}}(R_t) \equiv -\E\left( \frac{\partial^2 \ell}{\partial R_t^2} \right) = \frac{\E(I_t)}{R_t^2} = \frac{R_t \Lambda_t}{R_t^2} = \frac{\Lambda_t}{R_t}.
  \end{equation}
  \item 相对 Fisher 信息量（即无量纲标准化信息）：
  \begin{equation}
    \mathcal{F}(R_t) \equiv R_t^2 \, \mathcal{I}_{\text{Poisson}}(R_t) = R_t \Lambda_t = \E(I_t).
  \end{equation}
  这给出了极具穿透力的经典结论：\textbf{再生数的信息量在数值上恰好等于当期累积的期望病例数}。
\end{enumerate}

\subsection{四项核心数学连接与区别定理}

\begin{theorembox}{推导 14.1：向负二项超传播与环境随机性拓展的四大理论定理}
将本文的微观负二项分支与双目标框架与 Parag 的泊松框架对齐，可严格证明以下四项数学结论：
\end{theorembox}

\begin{enumerate}[label=\textbf{结论 \arabic*：}]
  \item \textbf{渐近退化性：无超传播且无环境波动时严格重合。}
  在本文模型中，令超传播参数 $k \to \infty$（无超级传播）、报告率 $\rho = 1$（全量报告）、祖代规模记为 $\Lambda_t$。根据式~\eqref{eq:fisher-def}：
  \begin{equation}
    \lim_{k \to \infty, \rho \to 1} \mathcal{I}(R_t) = \lim_{k \to \infty, \rho \to 1} \frac{n}{\frac{R_t}{\rho} + \frac{R_t^2}{k}} = \frac{n}{R_t + 0} = \frac{\Lambda_t}{R_t} = \mathcal{I}_{\text{Poisson}}(R_t).
  \end{equation}
  并且由于 $v_R = 0$，目标 1 与目标 2 重合，$A \to \Rbar = R_t$，$V(m) = 0 + \frac{R_t}{n} = \frac{1}{\mathcal{I}_{\text{Poisson}}}$。本文结论完整包含了 Parag 的经典结论作为特例。

  \item \textbf{超传播缓冲效应：不完全报告在负二项下的信息衰减小于泊松情形。}
  在泊松模型下，若报告率从 $1$ 降为 $\rho$，观测方差从 $R_t/n$ 升为 $R_t/(\rho n)$，相对信息损失比例严格等于 $1 - \rho$：
  \begin{equation}
    \frac{\mathcal{I}_{\rho=1} - \mathcal{I}_\rho}{\mathcal{I}_{\rho=1}} = \frac{\frac{n}{R_t} - \frac{\rho n}{R_t}}{\frac{n}{R_t}} = 1 - \rho.
  \end{equation}
  但在超传播（$k < \infty$）的真实世界下，代入式~\eqref{eq:fisher-def}：
  \begin{align}
    \frac{\mathcal{I}_{\rho=1} - \mathcal{I}_\rho}{\mathcal{I}_{\rho=1}} &= 1 - \frac{\frac{n}{R_t/\rho + R_t^2/k}}{\frac{n}{R_t + R_t^2/k}} = 1 - \frac{R_t + \frac{R_t^2}{k}}{\frac{R_t}{\rho} + \frac{R_t^2}{k}} \nonumber \\
    &= \frac{\left(\frac{R_t}{\rho} + \frac{R_t^2}{k}\right) - \left(R_t + \frac{R_t^2}{k}\right)}{\frac{R_t}{\rho} + \frac{R_t^2}{k}} = \frac{R_t\left(\frac{1}{\rho} - 1\right)}{\frac{R_t}{\rho}\left(1 + \frac{\rho R_t}{k}\right)} \nonumber \\
    &= \frac{1 - \rho}{1 + \frac{\rho R_t}{k}} < 1 - \rho.
  \end{align}
  由于分母中存在严格正项 $\frac{\rho R_t}{k} > 0$，超传播使得信息损失比例严格低于 $1 - \rho$！
  \textbf{直观解释：}超级传播个体造成的传染力方差本身已占主导，后置的二项抽样漏报只是在已经严重离散的数据上追加了一层次要随机性，因而边际信息破坏力相对减弱。

  \item \textbf{方差稳定变换在负二项下不可行。}
  在泊松模型中，Parag 利用 Anscombe 平方根变换 $Z_t = 2\sqrt{I_t}$，其渐近方差由一阶泰勒展开给出：
  \begin{equation}
    \Var(2\sqrt{I_t}) \approx \left[ 2 \cdot \frac{1}{2\sqrt{\E I_t}} \right]^2 \Var(I_t) = \frac{1}{R_t \Lambda_t} \cdot (R_t \Lambda_t) \equiv 1.
  \end{equation}
  由此彻底消除了方差对未知参数 $R_t$ 的依赖性。
  但在负二项超传播下，方差为 $\E(C_t) + \frac{\rho^2 n R_t^2}{k} = m R_t + \frac{\rho m R_t^2}{k}$。由于含有均值的二次项，任何幂函数变换均无法消去 $R_t$。

  \item \textbf{统计信息与物理离散度的根本边界。}
  Parag 的信息论体系是“\textbf{推断视角（Inferential perspective）}”：它回答“给定确定的现实，怎样从有限观测中最大限度抽取信息”。因此只要样本无限大，后验分布必然无限坍缩为狄拉克函数（Dirac delta），方差必然为零。
  而本文的环境随机性是“\textbf{生成视角（Generative perspective）}”：环境方差 $v_R$ 是大自然在每一世代施加的物理扰动。即便观测没有任何噪声且样本无限大，明天的传播强度相对于过去的均值依然具有物理发散度 $v_R$。这是生态学生物学现实，无法仅通过改进算法或堆积病例来消除。
\end{enumerate}
